% sec_entropy_asymptotics.tex -- Poisson smoothing and entropy asymptotics
\section{Poisson Smoothing and Entropy Asymptotics}%
\label{sec:precision-potential}

\paragraph{Assumption and dependency map.}
Throughout this section we write
$\Delta:=\gamma-\bar\gamma$ and
$\mu_n:=\int_{\RR}\Delta^n\,\dd\nu(\gamma)$ whenever the corresponding
moment exists.  For ease of reference, the moment hypotheses and logical
dependencies are as follows.
{\renewcommand{\arraystretch}{1.15}
\begin{center}
\small
\begin{tabular}{p{0.31\textwidth}p{0.23\textwidth}p{0.34\textwidth}}
\textbf{Result} & \textbf{Moment input} & \textbf{Main ingredients} \\
Proposition~\ref{thm:poisson-second-order}
& mean, variance, finite centred fourth moment
& direct Taylor expansion of $R_\varepsilon$ \\
Theorem~\ref{thm:poisson-kl-even-orders}
& finite centred $(2N+1)$st absolute moment
& Proposition~\ref{prop:cayley-mode-remainder} and
Lemma~\ref{lem:mode-parity-vanish} \\
Theorem~\ref{thm:poisson-kl-eighth}
& mean, variance, finite centred seventh absolute moment
& Proposition~\ref{prop:cayley-mode-remainder},
Lemma~\ref{lem:mode-parity-vanish}, and
Appendix~\ref{app:entropy-integrals} \\
Theorems~\ref{thm:poisson-defect-ladder}
and~\ref{thm:poisson-defect-amplification}
& mean, variance, finite centred seventh absolute moment
& the explicit coefficients from
Theorem~\ref{thm:poisson-kl-eighth} \\
Theorem~\ref{thm:poisson-defect-stability}
& mean, variance, finite centred fourth moment
& Lemma~\ref{lem:delta6-fourth-moment} (fourth-moment formula for $\Delta_6$)
\end{tabular}
\end{center}}
In particular, Appendix~\ref{app:entropy-integrals} enters only in the
explicit evaluation of the $t^{-6}$ and $t^{-8}$ coefficients in
Theorem~\ref{thm:poisson-kl-eighth};
Appendix~\ref{app:coefficient-assembly} provides the full step-by-step
coefficient assembly for verification.
The later defect identities are algebraic consequences of those
coefficients.

\subsection{Entropy in Cayley coordinates}

The logarithmic Jacobian of the Cayley coordinate is
\begin{equation}\label{eq:precision-potential}
\mathcal P(x):=-\log J(x),\qquad
J(x):=\upsilon'(x)=\frac{1}{\pi(1+x^{2})},
\end{equation}
equivalently,
\begin{equation}\label{eq:precision-potential-expanded}
\mathcal P(x)=\log\pi+\log(1+x^{2}).
\end{equation}

\begin{proposition}[Entropy identity in Cayley coordinates]%
\label{thm:precision-ledger}
Let $X$ be a real-valued random variable with density $f_X$ such that
$h_{\mathrm{diff}}(X)>-\infty$ and $\EE[\mathcal P(X)]<\infty$.  Set
$U:=\upsilon(X)\in\UU$.  Then
\begin{equation}\label{eq:entropy-ledger}
h_{\mathrm{diff}}(X)
=h_{\mathrm{diff}}(U)+\EE[\mathcal P(X)].
\end{equation}
\end{proposition}

\begin{proof}
The density of $U$ is
\[
f_U(u)=f_X(\tan(\pi u))\,\pi(1+\tan^2(\pi u)).
\]
The standard change-of-variables formula for differential entropy gives
\[
h_{\mathrm{diff}}(X)
=h_{\mathrm{diff}}(U)+\EE\!\left[\log\frac{1}{\abs{\upsilon'(X)}}\right]
=h_{\mathrm{diff}}(U)+\EE[\mathcal P(X)].
\]
\end{proof}

\begin{lemma}[Remainder bookkeeping]\label{lem:remainder-bookkeeping}
Let $\nu$ satisfy the assumptions of Lemma~\ref{lem:delta-expansion}
and let $M,N\ge 2$ be integers.
If $\delta_t$ is replaced by the truncation
\eqref{eq:delta-t-mode-expansion} with parameter $M$ and the Taylor
series of $\Phi$ is truncated after the $N$th term, then there exist
constants $C_{M}$ and $C_{M,N}$ (depending only on $M$, $N$, and the
moments of~$\nu$ up to order $M+1$) such that:
\begin{enumerate}[label=(\roman*)]
\item the remainder of the Taylor series contributes at most
      $C_{M,N}t^{-2(N+1)}$ in $L^1(\omega)$;
\item whenever one replaces $k$ of the factors of $\delta_t$ by the
      remainder $r_{M,t}$, the resulting $L^1(\omega)$ contribution is
      bounded by $C_{M} t^{-(M+1)} t^{-2(k-1)}$.
\end{enumerate}
Consequently, to control the coefficient of $t^{-2L}$ it is enough to
take $M=2L$ and $N=L$, in which case both sources of error contribute
$O(t^{-(2L+2)})$ in $L^1(\omega)$.
\end{lemma}

\begin{proof}
Part (i) is exactly Lemma~\ref{lem:delta-expansion}(ii)--(iii): the tail
of the Taylor series produces an $L^1(\omega)$ error bounded by
$C_{M,N}\|\delta_t\|_{L^\infty}^{N+1}=O(t^{-2(N+1)})$.
For (ii), each factor $r_{M,t}$ has $L^\infty$ norm $O(t^{-(M+1)})$
by Lemma~\ref{lem:delta-expansion}(i), while every remaining factor
of~$\delta_t$ contributes $O(t^{-2})$ in sup norm.
Integrating against the finite measure $\omega$ therefore yields the
$L^1(\omega)$ bound $C_{M} t^{-(M+1)} t^{-2(k-1)}$.
Choosing $M=2L$ and $N=L$ makes both error sources $O(t^{-(2L+2)})$,
which is the remainder claimed at the end of the lemma.
\end{proof}

\begin{lemma}[Entropy remainder order]\label{lem:entropy-remainder-order}
Let $L\ge 2$ and assume $\nu$ has a finite centred $(2L+1)$st absolute
moment.  Define
\[
\delta_t^{(2L)}(y):=\sum_{n=2}^{2L}\frac{\mu_n}{t^n}u_n(y),
\qquad y\in\RR,
\]
and keep $\Phi(z)=(1+z)\log(1+z)-z$.  Then, uniformly in $y$,
\[
\delta_t(y)=\delta_t^{(2L)}(y)+O(t^{-(2L+1)}),
\]
and in $L^1(\omega)$,
\[
\Phi(\delta_t)
=\sum_{k=2}^{L}\frac{(-1)^k}{k(k-1)}\bigl(\delta_t^{(2L)}\bigr)^k
  +O(t^{-(2L+2)}).
\]
Consequently,
\[
D_{\mathrm{KL}}(h_t\|g_t)
=\sum_{k=2}^{L}\frac{(-1)^k}{k(k-1)}
  \int_{\RR}\frac{\bigl(\delta_t^{(2L)}(y)\bigr)^k}{\pi(1+y^2)}\dd y
  +O(t^{-(2L+2)}).
\]
\end{lemma}

\begin{proof}
Lemma~\ref{lem:delta-expansion}(i) with $M=2L$ gives the uniform bound
on $\delta_t-\delta_t^{(2L)}$.  Applying
Lemma~\ref{lem:remainder-bookkeeping} with the same choice of $M$ and
$N=L$ shows that truncating the Taylor series of $\Phi$ introduces an
$L^1(\omega)$ error of size $O(t^{-(2L+2)})$.  Integrating against
$\omega(\dd y)=[\pi(1+y^2)]^{-1}\dd y$ yields the final identity.
\end{proof}

\begin{remark}
The preceding proposition is a standard change-of-variables formula for
differential entropy.  It is included to fix notation and to make the
role of the Cayley Jacobian explicit for the subsequent mode calculus.
\end{remark}

\subsection{Poisson smoothing}

For $t>0$ and $x\in\RR$ define
\begin{equation}\label{eq:poisson-kernel}
P_t(x):=\frac{1}{\pi}\frac{t}{x^{2}+t^{2}}.
\end{equation}

\begin{proposition}[Basic identities for the Poisson kernels]%
\label{prop:poisson-identities}
For $s,t>0$ one has
\begin{equation}\label{eq:poisson-basic}
P_t(x)=\frac{1}{t}P_1\!\left(\frac{x}{t}\right),
\qquad
\widehat{P_t}(\xi)=e^{-t\abs{\xi}},
\qquad
P_t*P_s=P_{t+s}.
\end{equation}
In particular,
\[
\mu=P_1(x)\dd x,
\qquad
P_t*\mu=P_{t+1}.
\]
\end{proposition}

\begin{proof}
The scaling identity is immediate from the definition.  The Fourier
transform of~$P_1$ is $e^{-\abs{\xi}}$
\cite[Chapter~III]{Stein1971}; scaling yields
$\widehat{P_t}(\xi)=e^{-t\abs{\xi}}$.  The convolution law follows by
multiplying Fourier transforms.
\end{proof}

\begin{proposition}[Standard facts on the Poisson semigroup]%
\label{thm:poisson-semigroup-law}
For $1\le p\le\infty$ and $t>0$ define
$T_t f:=P_t*f$.  Then:
\begin{enumerate}[label=(\roman*)]
\item $\|T_t f\|_{L^p(\RR)}\le \|f\|_{L^p(\RR)}$ for every $f\in L^p$;
\item $T_{t+s}=T_tT_s$ for all $s,t>0$;
\item for $1\le p<\infty$, $T_t f\to f$ in $L^p(\RR)$ as $t\downarrow 0$;
\item on $L^2(\RR)$, $\widehat{T_t f}(\xi)=e^{-t\abs{\xi}}\widehat f(\xi)$
      and the generator is $-|D|$.
\end{enumerate}
\end{proposition}

\begin{proof}
Since $\int_{\RR}P_t(x)\dd x=1$ and $P_t\ge 0$, Young's inequality gives
(i).  The semigroup law is Proposition~\ref{prop:poisson-identities}.
Strong continuity for $1\le p<\infty$ follows because $\{P_t\}_{t>0}$ is
an approximate identity.  The $L^2$-multiplier statement is exactly
\eqref{eq:poisson-basic}; the generator follows from
\[
\frac{e^{-t\abs{\xi}}-1}{t}\longrightarrow -\abs{\xi}
\qquad(t\downarrow 0)
\]
in $L^2$, which identifies the generator.
\end{proof}

\begin{remark}
The preceding proposition is standard and is included only to fix the
normalization of the Poisson kernel and of the generator $-|D|$ used in
the sequel.
\end{remark}

\subsection{Large-time comparison with a translated Poisson profile}

\begin{proposition}[Second-order approximation]%
\label{thm:poisson-second-order}
Let $\nu$ be a probability measure on~$\RR$ with mean
\[
\bar\gamma:=\int_{\RR}\gamma\dd\nu(\gamma),
\]
variance
\[
\sigma^{2}:=\int_{\RR}(\gamma-\bar\gamma)^{2}\dd\nu(\gamma),
\]
and finite centred fourth moment.  Define
\[
h_t:=P_t*\nu,
\qquad
g_t(x):=P_t(x-\bar\gamma).
\]
Let
\[
V(y):=\frac{3y^{2}-1}{(1+y^{2})^{2}}.
\]
Then there exists a constant $C_\nu>0$ such that for every $t\ge 1$,
\begin{equation}\label{eq:poisson-profile}
\sup_{x\in\RR}\left|
\frac{h_t(x)}{g_t(x)}
-1
-\frac{\sigma^{2}}{t^{2}}V\!\left(\frac{x-\bar\gamma}{t}\right)
\right|
\le \frac{C_\nu}{t^{3}}.
\end{equation}
Consequently,
\begin{equation}\label{eq:poisson-l1-rate}
\|h_t-g_t\|_{L^{1}(\RR)}=O(t^{-2}),
\qquad
\|h_t-g_t\|_{L^\infty(\RR)}=O(t^{-3}).
\end{equation}
Moreover,
\begin{equation}\label{eq:poisson-kl-rate}
D_{\mathrm{KL}}(h_t\|g_t)=O(t^{-4}).
\end{equation}
\end{proposition}

\begin{proof}
Write $\Delta:=\gamma-\bar\gamma$ and $x=\bar\gamma+ty$.  Then
\[
\frac{P_t(x-\gamma)}{g_t(x)}
=\frac{1+y^{2}}{1+(y-\Delta/t)^{2}}.
\]
Set $\varepsilon:=\Delta/t$.  A direct algebraic computation gives the
exact identity
\begin{align*}
\frac{1+y^{2}}{1+(y-\varepsilon)^{2}}
&=1+\frac{2y}{1+y^{2}}\varepsilon
+\frac{3y^{2}-1}{(1+y^{2})^{2}}\varepsilon^{2} \\
&\quad
+\frac{4y(y^{2}-1)\varepsilon^{3}-(3y^{2}-1)\varepsilon^{4}}
{(1+y^{2})^{2}(1+(y-\varepsilon)^{2})}.
\end{align*}
Since
\[
\frac{4\abs{y}\abs{y^{2}-1}}{(1+y^{2})^{2}}\le 2,
\qquad
\frac{\abs{3y^{2}-1}}{(1+y^{2})^{2}}\le 1,
\]
the remainder is bounded by $2\abs{\varepsilon}^{3}+\abs{\varepsilon}^{4}$
uniformly in~$y$.  Taking expectation and using $\EE[\Delta]=0$ yields
\[
\frac{h_t(\bar\gamma+ty)}{g_t(\bar\gamma+ty)}
=1+\frac{\sigma^{2}}{t^{2}}V(y)+R_t(y),
\]
with
\[
\sup_{y\in\RR}\abs{R_t(y)}
\le \frac{2\EE[\abs{\Delta}^{3}]+\EE[\abs{\Delta}^{4}]}{t^{3}}
\]
for $t\ge 1$.  This proves \eqref{eq:poisson-profile}.

The function~$U$ is bounded, so
$\sup_x\abs{h_t(x)/g_t(x)-1}=O(t^{-2})$.  Since $\int_{\RR}g_t=1$,
\[
\|h_t-g_t\|_{L^1}
\le \sup_{x}\abs{\frac{h_t(x)}{g_t(x)}-1}\int_{\RR}g_t(x)\dd x
=O(t^{-2}).
\]
Also $\|g_t\|_{L^\infty}=1/(\pi t)$, so
\[
\|h_t-g_t\|_{L^\infty}
\le \|g_t\|_{L^\infty}\sup_{x}\abs{\frac{h_t(x)}{g_t(x)}-1}
=O(t^{-3}).
\]
Let
\[
\delta_t(x):=\frac{h_t(x)}{g_t(x)}-1.
\]
The bound \eqref{eq:poisson-profile} implies that
$\|\delta_t\|_{L^\infty}=O(t^{-2})$, so for large $t$ one has
$\|\delta_t\|_{L^\infty}\le 1/2$.  Since
$\int g_t\delta_t\dd x=\int(h_t-g_t)\dd x=0$ and
\[
(1+s)\log(1+s)-s\le 2s^{2}
\qquad(\abs{s}\le 1/2),
\]
it follows that
\[
D_{\mathrm{KL}}(h_t\|g_t)
=\int_{\RR}g_t(x)\bigl((1+\delta_t(x))\log(1+\delta_t(x))-\delta_t(x)\bigr)\dd x
\le 2\|\delta_t\|_{L^\infty}^{2}.
\]
Hence \eqref{eq:poisson-kl-rate} holds.
\end{proof}

\subsection{Cayley--Chebyshev mode calculus}

For $\varepsilon\in\RR$ and $y\in\RR$ set
\begin{equation}\label{eq:cayley-comparison-kernel}
R_\varepsilon(y):=\frac{1+y^2}{1+(y-\varepsilon)^2}.
\end{equation}
For fixed $y$, the Taylor expansion of $R_\varepsilon(y)$ at
$\varepsilon=0$ defines coefficients $u_n(y)$ through
\begin{equation}\label{eq:cayley-mode-series}
R_\varepsilon(y)=\sum_{n=0}^{\infty}u_n(y)\varepsilon^n,
\qquad
|\varepsilon|<\sqrt{1+y^2}.
\end{equation}

\begin{theorem}[Cayley--Chebyshev mode formula]%
\label{thm:cayley-chebyshev-mode}
Let $U_n$ denote the Chebyshev polynomials of the second kind.  If
$y=\tan\theta$ with $\theta\in(-\pi/2,\pi/2)$, then for every $n\ge 0$,
\begin{equation}\label{eq:cayley-chebyshev-mode}
u_n(\tan\theta)=\cos^n\theta\,U_n(\sin\theta).
\end{equation}
For $n\ge 1$ this is equivalent to
\begin{equation}\label{eq:cayley-chebyshev-trig}
u_n(\tan\theta)
=\cos^{n-1}\theta\,
\sin\!\Bigl((n+1)\Bigl(\frac{\pi}{2}-\theta\Bigr)\Bigr).
\end{equation}
Consequently
\begin{equation}\label{eq:cayley-mode-parity}
u_n(-y)=(-1)^n u_n(y)
\qquad(n\ge 0).
\end{equation}
Moreover, for every integer $m\ge 1$,
\begin{equation}\label{eq:cayley-mode-even-fourier}
u_{2m}(\tan\theta)
=(-1)^m 2^{-(2m-1)}
\sum_{j=0}^{2m-1}\binom{2m-1}{j}\cos(2(j+1)\theta),
\end{equation}
and for every integer $m\ge 0$,
\begin{equation}\label{eq:cayley-mode-odd-fourier}
u_{2m+1}(\tan\theta)
=(-1)^m 2^{-2m}
\sum_{j=0}^{2m}\binom{2m}{j}\sin(2(j+1)\theta).
\end{equation}
\end{theorem}

\begin{proof}
With $y=\tan\theta$ one has
\[
R_\varepsilon(\tan\theta)
=\frac{1}{1-2(\varepsilon\cos\theta)\sin\theta+(\varepsilon\cos\theta)^2}.
\]
The generating function
\[
\sum_{n=0}^{\infty}U_n(x)r^n=\frac{1}{1-2xr+r^2}
\]
therefore yields \eqref{eq:cayley-chebyshev-mode} with
$x=\sin\theta$ and $r=\varepsilon\cos\theta$.  Since
\[
U_n(\sin\theta)
=\frac{\sin((n+1)(\pi/2-\theta))}{\cos\theta},
\]
formula \eqref{eq:cayley-chebyshev-trig} follows for $n\ge 1$.  The
parity relation \eqref{eq:cayley-mode-parity} is immediate from
\eqref{eq:cayley-chebyshev-trig}.

For the finite Fourier expansions, write
\[
u_{2m}(\tan\theta)=(-1)^m\cos^{2m-1}\theta\,\cos((2m+1)\theta),
\]
\[
u_{2m+1}(\tan\theta)=(-1)^m\cos^{2m}\theta\,\sin((2m+2)\theta),
\]
and use
\[
1+e^{2\ii\theta}=2e^{\ii\theta}\cos\theta.
\]
Taking real and imaginary parts after the binomial expansion gives
\eqref{eq:cayley-mode-even-fourier} and
\eqref{eq:cayley-mode-odd-fourier}.
\end{proof}

\begin{proposition}[Finite mode truncations]%
\label{prop:cayley-mode-remainder}
For every integer $N\ge 0$ there exists a rational function
$v_N:\RR^2\to\RR$, bounded on~$\RR^2$, such that
\begin{equation}\label{eq:cayley-mode-remainder}
R_\varepsilon(y)
=\sum_{n=0}^{N}u_n(y)\varepsilon^n+\varepsilon^{N+1}v_N(y,\varepsilon)
\qquad(y,\varepsilon\in\RR).
\end{equation}
\end{proposition}

\begin{proof}
For every $x$ and $r$ one has the finite-sum identity
\[
\sum_{n=0}^{N}U_n(x)r^n
=\frac{1-r^{N+1}U_{N+1}(x)+r^{N+2}U_N(x)}{1-2xr+r^2}.
\]
Substituting $x=\sin\theta$ and $r=\varepsilon\cos\theta$ gives
\eqref{eq:cayley-mode-remainder} with
\[
v_N(\tan\theta,\varepsilon)
=\frac{\cos^{N+1}\theta\,U_{N+1}(\sin\theta)
-\varepsilon\cos^{N+2}\theta\,U_N(\sin\theta)}
{1-2(\varepsilon\cos\theta)\sin\theta+(\varepsilon\cos\theta)^2}.
\]
This expression is rational in $(y,\varepsilon)$.  Since the denominator
is positive on $\RR^2$ and the numerator has lower degree than the
denominator, $v_N$ is bounded on~$\RR^2$.
\end{proof}

For later use we record the first modes explicitly:
\begin{equation}\label{eq:cayley-first-modes}
u_2(y)=\frac{3y^2-1}{(1+y^2)^2},
\qquad
u_3(y)=\frac{4y(y^2-1)}{(1+y^2)^3},
\qquad
u_4(y)=\frac{5y^4-10y^2+1}{(1+y^2)^4},
\end{equation}
\begin{equation}\label{eq:cayley-next-modes}
u_5(y)=\frac{6y^5-20y^3+6y}{(1+y^2)^5},
\qquad
u_6(y)=\frac{7y^6-35y^4+21y^2-1}{(1+y^2)^6}.
\end{equation}

\begin{lemma}[Mode parity cancellation]\label{lem:mode-parity-vanish}
For integers $n_1,\dots,n_k\ge 1$ let
\[
I_{n_1,\dots,n_k}:=
\int_{\RR}
\frac{u_{n_1}(y)\cdots u_{n_k}(y)}{\pi(1+y^2)}\,\dd y.
\]
Then $I_{n_1,\dots,n_k}=0$ whenever $n_1+\cdots+n_k$ is odd.
\end{lemma}

\begin{proof}
From \eqref{eq:cayley-mode-parity}, $u_n(-y)=(-1)^n u_n(y)$ for every
$n\ge1$, while $(1+y^2)^{-1}$ is even.  Hence the integrand is odd if
$n_1+\cdots+n_k$ is odd, so the integral vanishes.
\end{proof}

\subsection{Entropy asymptotics and rigidity}

\begin{lemma}[Asymptotic expansion of $\delta_t$ and the logarithmic nonlinearity]%
\label{lem:delta-expansion}
Let $\nu$ be a probability measure on $\RR$ with mean $\bar\gamma$ and
finite centred $(M+1)$st absolute moment, and let $M \ge 2$ be an integer.
With $\Delta:=\gamma-\bar\gamma$ and $\mu_n:=\int_\RR \Delta^n\,\dd\nu(\gamma)$, set
\[
\delta_t(\bar\gamma+ty)
:= \frac{h_t(\bar\gamma+ty)}{g_t(\bar\gamma+ty)}-1
= \int_\RR (R_{\Delta/t}(y)-1)\,\dd\nu(\gamma).
\]

\noindent\textbf{(i) Mode expansion of $\delta_t$.}
For every $M\ge 2$ and every $t\ge 1$,
\begin{equation}\label{eq:delta-t-mode-expansion}
\delta_t(\bar\gamma+ty)
= \sum_{n=2}^{M} \frac{\mu_n}{t^n} u_n(y) + r_{M,t}(y),
\end{equation}
where the remainder satisfies
\begin{equation}\label{eq:delta-t-remainder-bound}
\|r_{M,t}\|_{L^\infty(\RR)} \le \frac{C_M}{t^{M+1}},
\end{equation}
with $C_M:= \|v_M\|_{L^\infty}\int_\RR |\Delta|^{M+1}\,\dd\nu(\gamma)$
and $v_M$ the bounded rational function from
Proposition~\ref{prop:cayley-mode-remainder}.

\noindent\textbf{(ii) Taylor expansion of $\Phi$ with remainder.}
The function $\Phi(z):=(1+z)\log(1+z)-z$ has the convergent power series
\begin{equation}\label{eq:phi-power-series}
\Phi(z) = \sum_{k=2}^{\infty} \frac{(-1)^k}{k(k-1)}\,z^k, \qquad |z|<1,
\end{equation}
and for every integer $N\ge 2$ there exists a constant $C_N>0$ such that
\begin{equation}\label{eq:phi-remainder}
\left|\Phi(z) - \sum_{k=2}^{N}\frac{(-1)^k}{k(k-1)}\,z^k\right|
\le C_N\,|z|^{N+1}, \qquad |z|\le \tfrac12.
\end{equation}
One may take $C_N = 4/(N(N-1))$ explicitly.

\noindent\textbf{(iii) Integration in $L^1(\omega)$.}
Since $\|\delta_t\|_{L^\infty}=O(t^{-2})$, for $t$ sufficiently large
$|\delta_t(y)|\le 1/2$ for all $y$.  Applying~\eqref{eq:phi-remainder} with
$z=\delta_t(y)$ and using the mode expansion~\eqref{eq:delta-t-mode-expansion},
one may substitute the finite sum $\sum_{k=2}^N (-1)^k/(k(k-1))\,\delta_t^k$,
expand each power using~\eqref{eq:delta-t-mode-expansion}, and integrate
term by term against $\omega(\dd y) = [\pi(1+y^2)]^{-1}\dd y$.  The resulting
error from the $\Phi$-remainder is $O(t^{-2(N+1)})$ in $L^1(\omega)$, and
the error from $r_{M,t}$ contributes $O(t^{-(M+1-2)})$ per factor of $\delta_t$.
Choosing $M$ and $N$ large enough, the full remainder after integration is
$o(t^{-2N})$.  Term-by-term integration is justified by dominated convergence,
since $\omega$ is a finite measure and $\|\delta_t\|_{L^\infty}\to 0$.
\end{lemma}

\begin{proof}
\textbf{Part (i).}
By Proposition~\ref{prop:cayley-mode-remainder} with truncation index $M$,
\[
R_{\Delta/t}(y) - 1
= \sum_{n=1}^{M} u_n(y)\,\frac{\Delta^n}{t^n}
+ \frac{\Delta^{M+1}}{t^{M+1}}\,v_M(y,\Delta/t).
\]
Since $u_1$ is odd and $\int \Delta\,\dd\nu=0$, the $n=1$ term integrates to
zero.  Taking expectation against $\nu$,
\[
\delta_t(\bar\gamma+ty)
= \sum_{n=2}^{M}\frac{\mu_n}{t^n}u_n(y)
+ \frac{1}{t^{M+1}}\int_\RR \Delta^{M+1}v_M(y,\Delta/t)\,\dd\nu(\gamma).
\]
Since $\|v_M\|_{L^\infty(\RR^2)}<\infty$
(Proposition~\ref{prop:cayley-mode-remainder}),
the remainder is bounded in sup norm by $\|v_M\|_{L^\infty}t^{-(M+1)}\int|\Delta|^{M+1}\dd\nu$.

\textbf{Part (ii).}
For $|z|\le 1/2$ the series converges absolutely, so the partial-sum
remainder is bounded by
\[
\sum_{k=N+1}^\infty \frac{|z|^k}{k(k-1)}
\le \frac{|z|^{N+1}}{N(N-1)}\sum_{k=0}^\infty |z|^k
\le \frac{4}{N(N-1)}\,|z|^{N+1}.
\]

\textbf{Part (iii).}
The $L^1(\omega)$ error from the $\Phi$-truncation at order $N$ applied to
$\delta_t=O(t^{-2})$ is bounded by
\[
C_N\int_\RR|\delta_t(y)|^{N+1}\omega(\dd y)
\le C_N\|\delta_t\|_{L^\infty}^{N+1}=O(t^{-2(N+1)}).
\]
Each occurrence of $r_{M,t}$ in a product $\delta_t^{k-1}\cdot r_{M,t}$
contributes $O(t^{-2(k-1)}\cdot t^{-(M+1)})$ per factor.  For the
term-by-term integration, dominant convergence applies because $\omega$
is a probability measure and all integrands are bounded uniformly in
$y$ for each fixed $t$.
\end{proof}

\begin{theorem}[Odd orders vanish in the entropy expansion]%
\label{thm:poisson-kl-even-orders}
Let $N\ge 2$ and let $\nu$ be a probability measure on~$\RR$ with mean
$\bar\gamma$ and finite centred $(2N+1)$st absolute moment.  Set
\[
\Delta:=\gamma-\bar\gamma,
\qquad
\mu_n:=\int_{\RR}\Delta^n\dd\nu(\gamma),
\]
and keep the notation
\[
h_t:=P_t*\nu,
\qquad
g_t(x):=P_t(x-\bar\gamma).
\]
Then there exist real coefficients $c_4,c_6,\dots,c_{2N}$ such that
\begin{equation}\label{eq:poisson-kl-even-orders}
D_{\mathrm{KL}}(h_t\|g_t)
=\sum_{m=2}^{N}c_{2m}\,t^{-2m}+o(t^{-2N})
\qquad(t\to\infty).
\end{equation}
In particular, every odd power $t^{-(2m+1)}$ has zero coefficient.
\end{theorem}

\begin{proof}
Apply Lemma~\ref{lem:delta-expansion}(i) with $M=2N$: uniformly in $y$,
\[
\delta_t(\bar\gamma+ty)
=\sum_{n=2}^{2N}\frac{\mu_n}{t^n}u_n(y)+O(t^{-(2N+1)}).
\]
By Lemma~\ref{lem:delta-expansion}(ii)--(iii) and
Lemma~\ref{lem:entropy-remainder-order} (applied with $L=N$), one may
replace $\delta_t$ by its $2N$-term truncation, substitute the $N$-term
Taylor polynomial for $\Phi$, and integrate term by term in
$L^1(\omega)$, yielding
\[
D_{\mathrm{KL}}(h_t\|g_t)
=\int_{\RR}\frac{\Phi(\delta_t(\bar\gamma+ty))}{\pi(1+y^2)}\dd y
=\sum_{k=2}^{N}\frac{(-1)^k}{k(k-1)}
\sum_{\substack{n_1,\dots,n_k\ge 2\\n_1+\cdots+n_k=m}}
\frac{\mu_{n_1}\cdots\mu_{n_k}}{t^m}\,I_{n_1,\dots,n_k}
+o(t^{-2N}),
\]
where $I_{n_1,\dots,n_k}$ is defined in Lemma~\ref{lem:mode-parity-vanish}.
By that lemma all terms with $m$ odd vanish.  Reindexing $m=2r$ gives
\[
D_{\mathrm{KL}}(h_t\|g_t)=\sum_{r=2}^{N} c_{2r}\,t^{-2r}+o(t^{-2N}),
\]
which is \eqref{eq:poisson-kl-even-orders}.
\end{proof}

\begin{remark}[Moment input versus coefficient content]
\label{rem:even-order-moment-input}
The centred $(2N+1)$st absolute moment assumption in
Theorem~\ref{thm:poisson-kl-even-orders} is a remainder-control
hypothesis, not a claim that the coefficient of $t^{-2N}$ depends on
moments up to order $2N+1$.  Once the mode expansion is truncated at
order $2N$, the coefficient of $t^{-2N}$ is determined by moments up to
order $2N$; the extra absolute moment enters only through the uniform
bound on the remainder $r_{2N,t}$.  A more delicate truncation scheme
should reduce this hypothesis, but that refinement is not needed for
the present parity argument.
\end{remark}

\begin{theorem}[Eighth-order entropy expansion]%
\label{thm:poisson-kl-eighth}
Let $\nu$ be a probability measure on~$\RR$ with mean~$\bar\gamma$,
variance $\sigma^2>0$, and finite centred seventh absolute moment.  Then,
as $t\to\infty$,
\begin{equation}\label{eq:poisson-kl-eighth}
\begin{split}
D_{\mathrm{KL}}(h_t\|g_t)
=\frac{\sigma^4}{8t^4}
+\frac{\sigma^6+6\mu_3^2-8\sigma^2\mu_4}{64\,t^6}
\\
\quad
+\frac{3\sigma^8-12\sigma^4\mu_4+9\sigma^2\mu_3^2}{256\,t^8}
+\frac{12\sigma^2\mu_6-30\mu_3\mu_5+20\mu_4^2}{256\,t^8}
+o(t^{-8}).
\end{split}
\end{equation}
\end{theorem}

\begin{proof}
By Lemma~\ref{lem:delta-expansion}(i) with $M=6$,
\[
\delta_t(\bar\gamma+ty)
=\frac{\sigma^2}{t^2}u_2(y)
+\frac{\mu_3}{t^3}u_3(y)
+\frac{\mu_4}{t^4}u_4(y)
+\frac{\mu_5}{t^5}u_5(y)
+\frac{\mu_6}{t^6}u_6(y)
+r_{6,t}(y),
\]
uniformly in $y\in\RR$, with $\|r_{6,t}\|_{L^\infty}\le C_6 t^{-7}$.
By Lemma~\ref{lem:delta-expansion}(ii) with $N=4$ (since $\delta_t=O(t^{-2})$
and contributions from $k\ge 5$ are $O(t^{-10})$),
\[
\Phi(\delta_t)
=\frac12\delta_t^2-\frac16\delta_t^3+\frac1{12}\delta_t^4+O(t^{-10}).
\]
Also $g_t(\bar\gamma+ty)\dd(\bar\gamma+ty)=\pi^{-1}(1+y^2)^{-1}\dd y$.
Lemma~\ref{lem:entropy-remainder-order} with $L=4$ (equivalently,
$M=6$ and $N=4$ in Lemma~\ref{lem:remainder-bookkeeping}) justifies the
term-by-term integration against $\omega(\dd y)=[\pi(1+y^2)]^{-1}\dd y$:
the $\Phi$-tail and every occurrence of $r_{6,t}$ contribute
$o(t^{-8})$ in $L^1(\omega)$.
Write
\[
a_2:=\sigma^2u_2,\qquad a_3:=\mu_3u_3,\qquad a_4:=\mu_4u_4,
\qquad a_5:=\mu_5u_5,\qquad a_6:=\mu_6u_6,
\]
so that $\delta_t=t^{-2}a_2+t^{-3}a_3+t^{-4}a_4+t^{-5}a_5+t^{-6}a_6+O(t^{-7})$.
Up to order $t^{-8}$ one therefore has
\[
\frac12\delta_t^2
=\frac{a_2^2}{2t^4}
+\frac{a_2a_3}{t^5}
+\frac{a_2a_4+\frac12 a_3^2}{t^6}
+\frac{a_2a_5+a_3a_4}{t^7}
+\frac{a_2a_6+a_3a_5+\frac12 a_4^2}{t^8}
+O(t^{-9}),
\]
\[
-\frac16\delta_t^3
=-\frac{a_2^3}{6t^6}
-\frac{a_2^2a_3}{2t^7}
-\frac{\frac12 a_2^2a_4+\frac12 a_2a_3^2}{t^8}
+O(t^{-9}),
\]
By Lemma~\ref{lem:mode-parity-vanish}, the odd $t^{-5}$ and $t^{-7}$ terms in the first two expansions integrate to zero against $\pi^{-1}(1+y^2)^{-1}\dd y$.
\[
\frac1{12}\delta_t^4
=\frac{a_2^4}{12t^8}+O(t^{-9}).
\]
After integration against $\pi^{-1}(1+y^2)^{-1}\dd y$, the weighted
identities from Appendix~\ref{app:entropy-integrals} give
\[
\frac12\int_{\RR}\frac{u_2(y)^2}{\pi(1+y^2)}\dd y=\frac18,
\]
\[
\int_{\RR}\frac{\sigma^2\mu_4u_2u_4+\frac12\mu_3^2u_3^2-\frac16\sigma^6u_2^3}
{\pi(1+y^2)}\dd y
=\frac{\sigma^6+6\mu_3^2-8\sigma^2\mu_4}{64},
\]
and
\begin{align*}
\int_{\RR}
\frac{\sigma^2\mu_6u_2u_6+\mu_3\mu_5u_3u_5+\frac12\mu_4^2u_4^2}
{\pi(1+y^2)}\dd y
&=\frac{12\sigma^2\mu_6-30\mu_3\mu_5+20\mu_4^2}{256},
\\
\int_{\RR}
\frac{-\frac12\sigma^4\mu_4u_2^2u_4-\frac12\sigma^2\mu_3^2u_2u_3^2
+\frac1{12}\sigma^8u_2^4}{\pi(1+y^2)}\dd y
&=\frac{3\sigma^8-12\sigma^4\mu_4+9\sigma^2\mu_3^2}{256}.
\end{align*}
This proves \eqref{eq:poisson-kl-eighth}.  A step-by-step verification of
the coefficient assembly from the mode integrals is given in
Appendix~\ref{app:coefficient-assembly}.
\end{proof}

\begin{remark}[The seventh moment assumption is technical]
\label{rem:eighth-order-moment-input}
The coefficients in \eqref{eq:poisson-kl-eighth} involve only moments
through order six.  The finite centred seventh absolute moment
assumption is used only to obtain the uniform remainder
$r_{6,t}=O(t^{-7})$ in Lemma~\ref{lem:delta-expansion} and thereby
justify the $o(t^{-8})$ error after integration.  In that sense the
assumption is technical rather than conceptual.  We have not attempted
to optimize it here; a weaker hypothesis combining sixth-moment control
with a finer domination argument should suffice.
\end{remark}

\begin{remark}[Sanity check: symmetric two-point law]\label{rem:two-point-sanity}
For the symmetric two-point measure
$\nu=\tfrac12(\delta_{\bar\gamma+\sigma}+\delta_{\bar\gamma-\sigma})$,
the centered normalized variable satisfies $X=\pm 1$ with equal weight, so
$\beta_3=\mu_3/\sigma^3=0$, $\beta_4=\mu_4/\sigma^4=1$,
$\mu_5=0$, $\mu_6=\sigma^6$.
Substituting into the formula for $C_6$ yields
\[
C_6=\frac{\sigma^6+0-8\sigma^2\cdot\sigma^4}{64}
=\frac{\sigma^6-8\sigma^6}{64}=-\frac{7\sigma^6}{64}.
\]
For $C_8$, with $\mu_4=\sigma^4$, $\mu_5=0$, $\mu_6=\sigma^6$:
\[
C_8=\frac{3\sigma^8-12\sigma^4\cdot\sigma^4+0+12\sigma^2\cdot\sigma^6}{256}
+\frac{0+0+20\sigma^8}{256}
=\frac{3\sigma^8-12\sigma^8+12\sigma^8+20\sigma^8}{256}
=\frac{23\sigma^8}{256}.
\]
So the entropy expansion for the symmetric two-point law reads
\[
D_{\mathrm{KL}}(h_t\|g_t)=\frac{\sigma^4}{8t^4}-\frac{7\sigma^6}{64t^6}+\frac{23\sigma^8}{256t^8}+o(t^{-8}),
\]
with $C_6=-7\sigma^6/64$ and $C_8=23\sigma^8/256$.  Both defects
$\Delta_6=0$ and $\Delta_8=0$, consistently with
Theorems~\ref{thm:poisson-defect-ladder}
and~\ref{thm:poisson-defect-amplification}.
\end{remark}

\begin{remark}\label{rem:entropy-expansion-context}
The fourth-order coefficient $\sigma^4/8$ corresponds to the classical
Fisher-information rate for Cauchy smoothing.  At sixth order, the
coefficient $C_6$ already combines the skewness and excess kurtosis of
$\nu$, but its non-negativity reformulation is essentially Pearson's
classical inequality in adapted coordinates.  The genuinely new content
appears at eighth order: the coefficient $C_8$ brings in moments
through order six and is where the defect-ladder structure emerges.

In the entropic CLT literature
\cite{BobkovChistyakovGotzeEdgeworth2013}, Edgeworth-type corrections
for Gaussian targets are studied; the present expansion is for the
Poisson--Cauchy pair and uses the trigonometric mode calculus rather
than cumulant expansions.
For stable-law Fisher information and convergence rates,
Bobkov, Chistyakov, and G\"otze \cite{BobkovChistyakovGotzeFisher2014}
establish Fisher-information-based rates for convergence to stable laws.
For stable/Cauchy entropy dissipation, Johnson \cite{Johnson2013}
establishes a de Bruijn identity for symmetric stable laws, and
Toscani \cite{Toscani2015,Toscani2016} studies fractional Fisher
information and entropy inequalities.  Those results concern the
instantaneous entropy production rate along stable semigroup flows,
whereas the present contribution is a fixed-$\nu$, large-$t$
asymptotic \emph{expansion} of
$D_{\mathrm{KL}}(P_t*\nu\|P_t(\cdot-\bar\gamma))$ with explicit
$t^{-6}$ and $t^{-8}$ coefficients and defect-amplification
inequalities.  The two viewpoints are complementary: the de Bruijn
identity governs the leading-order derivative, while the present
expansion isolates higher-order structural invariants.
\end{remark}

\begin{center}
\fbox{\parbox{0.94\textwidth}{
\textbf{Notation guide for the defect formulas.}
Write
\[
C_6:=\frac{\sigma^6+6\mu_3^2-8\sigma^2\mu_4}{64}.
\]
\[
\begin{aligned}
C_8&:=\frac{3\sigma^8-12\sigma^4\mu_4+9\sigma^2\mu_3^2+12\sigma^2\mu_6}{256} \\
&\quad +\frac{-30\mu_3\mu_5+20\mu_4^2}{256}.
\end{aligned}
\]
so that Theorem~\ref{thm:poisson-kl-eighth} reads
\[
D_{\mathrm{KL}}(h_t\|g_t)=\frac{\sigma^4}{8t^4}+\frac{C_6}{t^6}+\frac{C_8}{t^8}+o(t^{-8}).
\]
The associated sixth- and eighth-order defects are
\[
\Delta_6(\nu):=-\frac{7}{64}\sigma^6-C_6,
\qquad
\Delta_8(\nu):=C_8-\frac{23}{256}\sigma^8.
\]
When only a finite fourth centred moment is assumed,
Lemma~\ref{lem:delta6-fourth-moment} provides the equivalent
fourth-moment formula \eqref{eq:delta6-fourth-moment-lemma} for $\Delta_6(\nu)$.
}}
\end{center}

\begin{theorem}[A two-level defect ladder]%
\label{thm:poisson-defect-ladder}
Let $\nu$ be a probability measure on~$\RR$ with mean~$\bar\gamma$,
variance $\sigma^2>0$, and finite centred seventh absolute moment.  Define
\begin{align*}
C_6&:=\frac{\sigma^6+6\mu_3^2-8\sigma^2\mu_4}{64}, \\
C_8&:=
\frac{3\sigma^8-12\sigma^4\mu_4+9\sigma^2\mu_3^2+12\sigma^2\mu_6}{256}
+\frac{-30\mu_3\mu_5+20\mu_4^2}{256}.
\end{align*}
Let $X:=(\gamma-\bar\gamma)/\sigma$ on the probability space
$(\RR,\nu)$, set
\[
\beta_n:=\EE[X^n],\qquad n\in\NN,\ n\ge1,\qquad
\beta_3:=\EE[X^3],
\qquad
Q_2:=X^2-1-\beta_3X,
\]
and define
\[
Q_3:=XQ_2-\|Q_2\|_{L^2(\nu)}^2X-\frac{\beta_3}{4}Q_2.
\]
This definition is purely algebraic: $Q_3$ is the linear combination of
$XQ_2$, $X$, and $Q_2$ that cancels the $X$ and $Q_2$ coefficients
emerging in \eqref{eq:defect-c8}, so no orthogonal projection in
$L^2(\nu)$ is being taken.
Then
\begin{equation}\label{eq:defect-c6}
\frac{64}{\sigma^6}C_6
=-7-2\beta_3^2-8\|Q_2\|_{L^2(\nu)}^2,
\end{equation}
and
\begin{equation}\label{eq:defect-c8}
\frac{256}{\sigma^8}C_8
=23
+12\|Q_3\|_{L^2(\nu)}^2
+32\|Q_2\|_{L^2(\nu)}^4
+\frac{61}{4}\beta_3^2\|Q_2\|_{L^2(\nu)}^2
+52\|Q_2\|_{L^2(\nu)}^2
+2\beta_3^4
+13\beta_3^2.
\end{equation}
In particular
\begin{equation}\label{eq:defect-c8-lower}
C_8\ge \frac{23}{256}\sigma^8,
\end{equation}
and equality holds if and only if
\[
\nu=\frac12\delta_{\bar\gamma+\sigma}
+\frac12\delta_{\bar\gamma-\sigma}.
\]
If $\tau:=\beta_3$ and $\kappa:=\beta_4$ denote the normalized
skewness and kurtosis, then
$\|Q_2\|_{L^2(\nu)}^2=\kappa-1-\tau^2$, so \eqref{eq:defect-c6} is
equivalent to the classical Pearson inequality
$\kappa-\tau^2\ge 1$
\cite{KlaassenVanEs2018,OgasawaraExtensions2017,SharmaBhandari2015}.
Identity \eqref{eq:defect-c8} is the
first higher-order component beyond Pearson's relation.
\end{theorem}

\begin{proof}
Write
\[
q_2:=\|Q_2\|_{L^2(\nu)}^2=\EE[Q_2^2],
\qquad
L:=\EE[XQ_2^2],
\qquad
M:=\EE[Q_2^3].
\]
Since $Q_2=X^2-1-\beta_3X$, one has
\[
q_2=\beta_4-1-\beta_3^2,
\]
and therefore
\[
\frac{64}{\sigma^6}C_6
=1+6\beta_3^2-8\beta_4
=-7-2\beta_3^2-8q_2,
\]
which is exactly \eqref{eq:defect-c6}.

The identity $X^2=1+\beta_3X+Q_2$ yields
\[
\beta_5=\beta_3(\beta_3^2+2+2q_2)+L,
\]
\[
\beta_6=1+3\beta_3^2+3q_2+\beta_3^4+3\beta_3^2q_2+3\beta_3L+M.
\]
Hence, if
\[
P_8:=3-12\beta_4+9\beta_3^2+12\beta_6-30\beta_3\beta_5+20\beta_4^2,
\]
then a direct substitution gives
\[
P_8
=23+20q_2^2+16\beta_3^2q_2+64q_2+2\beta_3^4+13\beta_3^2+6\beta_3L+12M.
\]
On the other hand,
\[
\|Q_3\|_{L^2(\nu)}^2
=M+\frac{\beta_3}{2}L+q_2-q_2^2+\frac{\beta_3^2}{16}q_2.
\]
Substituting this identity into the preceding expression for $P_8$
yields \eqref{eq:defect-c8}.  Since
$P_8=256C_8/\sigma^8$, inequality \eqref{eq:defect-c8-lower} follows.

If equality holds in \eqref{eq:defect-c8-lower}, then every non-negative
term on the right-hand side of \eqref{eq:defect-c8} vanishes.  In
particular $\beta_3=0$ and $\|Q_2\|_{L^2(\nu)}=0$, so $X^2=1$
$\nu$-almost everywhere.  Since $\EE[X]=0$ and $\EE[X^2]=1$, it follows
that $X=\pm 1$ with equal weights.  Thus
\[
\nu=\frac12\delta_{\bar\gamma+\sigma}
+\frac12\delta_{\bar\gamma-\sigma}.
\]
The converse is immediate.
\end{proof}

\subsection{Algebraic meaning of the defect ladder}
\label{subsec:defect-ladder-geometry}

The coordinates $Q_2$ and $Q_3$ in Theorem~\ref{thm:poisson-defect-ladder}
are designed to isolate the moment combinations that survive in the
entropy coefficients.
They are defined only through algebraic cancellations, and we do not
interpret them as orthogonal projection residuals in $L^2(\nu)$.

Let $X:=(\gamma-\bar\gamma)/\sigma$ with $\EE[X]=0$ and $\EE[X^2]=1$.
The quadratic coordinate
\[
Q_2 := X^2 - 1 - \beta_3 X
\]
removes the components of $X^2-1$ in the span of $\{1,X\}$, so
$\|Q_2\|_{L^2(\nu)}^2=\beta_4-1-\beta_3^2$ reproduces the Pearson
skewness--kurtosis gap.

The cubic coordinate
\[
Q_3 := X Q_2 - \|Q_2\|_{L^2(\nu)}^2 X - \frac{\beta_3}{4} Q_2
\]
is chosen to cancel the lower-order interactions that appear when the
terms of Theorem~\ref{thm:poisson-kl-eighth} are regrouped.
It is not an orthogonal projection residual; instead it is the unique
linear combination of $X Q_2$, $X$, and $Q_2$ that removes the $X$ and
$Q_2$ coefficients in \eqref{eq:defect-c8}.
Consequently $\|Q_3\|_{L^2(\nu)}^2$ measures the remaining cubic
contribution after those cancellations.

With these coordinates, the formulas in Theorem~\ref{thm:poisson-defect-ladder}
show that:
\begin{itemize}
\item The $t^{-6}$ coefficient $C_6$ depends only on $\beta_3$ and
$\|Q_2\|_{L^2(\nu)}^2$, so the sixth-order defect is a combination of
the skewness and the Pearson gap.
\item The $t^{-8}$ coefficient $C_8$ adds the single new coordinate
$\|Q_3\|_{L^2(\nu)}^2$ together with polynomial combinations of
$\beta_3$ and $\|Q_2\|_{L^2(\nu)}^2$.
\item The symmetric two-point law
$\tfrac12(\delta_{\bar\gamma+\sigma}+\delta_{\bar\gamma-\sigma})$
forces $\beta_3=\|Q_2\|_{L^2(\nu)}=\|Q_3\|_{L^2(\nu)}=0$, so all defect
coordinates vanish simultaneously.
\end{itemize}

Thus the defect ladder provides an algebraic gauge of the departure
from the symmetric two-point extremizer, expressed through one quadratic
coordinate and one higher-order correction.

\begin{theorem}[Defect amplification]\label{thm:poisson-defect-amplification}
Let $\nu$ be a probability measure on~$\RR$ with mean~$\bar\gamma$,
variance $\sigma^2>0$, and finite centred seventh absolute moment.  Define
\[
\Delta_6(\nu):=-\frac{7}{64}\sigma^6-C_6,
\qquad
\Delta_8(\nu):=C_8-\frac{23}{256}\sigma^8,
\]
and set
\[
q_2:=\|Q_2\|_{L^2(\nu)}^2.
\]
Then
\[
\Delta_6(\nu)
=\frac{\sigma^6}{32}\beta_3^2+\frac{\sigma^6}{8}q_2,
\]
and
\begin{align}
\Delta_8(\nu)
&=
\frac{13\sigma^2}{8}\Delta_6(\nu)
\notag\\
&\quad
+\frac{\sigma^8}{256}\left(
12\|Q_3\|_{L^2(\nu)}^2
+32q_2^2
+\frac{61}{4}\beta_3^2q_2
+2\beta_3^4
\right). \label{eq:defect-amplification}
\end{align}
In particular,
\[
\Delta_8(\nu)\ge \frac{13\sigma^2}{8}\Delta_6(\nu),
\]
and equality holds if and only if
\[
\nu=\frac12\delta_{\bar\gamma+\sigma}
+\frac12\delta_{\bar\gamma-\sigma}.
\]
\end{theorem}

\begin{proof}
Equation~\eqref{eq:defect-c6} gives
\[
\Delta_6(\nu)
=-\frac{\sigma^6}{64}\left(7+\frac{64}{\sigma^6}C_6\right)
=\frac{\sigma^6}{32}\beta_3^2+\frac{\sigma^6}{8}q_2.
\]
Likewise, subtracting $23$ from \eqref{eq:defect-c8} yields
\[
\Delta_8(\nu)
=\frac{\sigma^8}{256}\Bigl(
12\|Q_3\|_{L^2(\nu)}^2
+32q_2^2
+\frac{61}{4}\beta_3^2q_2
+52q_2
+2\beta_3^4
+13\beta_3^2
\Bigr).
\]
By the preceding formula for $\Delta_6(\nu)$,
\[
\frac{13\sigma^2}{8}\Delta_6(\nu)
=\frac{\sigma^8}{256}\bigl(52q_2+13\beta_3^2\bigr),
\]
and substituting this identity proves \eqref{eq:defect-amplification}.

If equality holds in the lower bound, then every non-negative term in
\eqref{eq:defect-amplification} vanishes.  In particular,
$\beta_3=0$ and $q_2=0$, hence $Q_2=X^2-1$ vanishes
$\nu$-almost everywhere.  Since $\EE[X]=0$ and $\EE[X^2]=1$, it follows
that $X=\pm 1$ with equal weights, which is exactly the stated
two-point law.  The converse is immediate.
\end{proof}

\begin{remark}
The amplification inequality
$\Delta_8\ge(13\sigma^2/8)\Delta_6$ says that any departure from the
symmetric two-point law measured by the sixth-order defect is strictly
magnified at eighth order, with the amplification factor
$13\sigma^2/8$ depending only on the variance.  This phenomenon is
invisible at the level of the fourth-order coefficient alone.
\end{remark}

\begin{lemma}[The sixth-order defect under finite fourth moment]%
\label{lem:delta6-fourth-moment}
Let $\nu$ be a probability measure on~$\RR$ with mean~$\bar\gamma$,
variance $\sigma^2>0$, and finite centred fourth moment.  Set
$X:=(\gamma-\bar\gamma)/\sigma$, $\beta_3:=\EE[X^3]$, $\beta_4:=\EE[X^4]$.
Then
\begin{equation}\label{eq:delta6-fourth-moment-lemma}
\Delta_6(\nu)
=\frac{\sigma^6}{32}\beta_3^2
+\frac{\sigma^6}{8}(\beta_4-1-\beta_3^2).
\end{equation}
In particular, $\Delta_6(\nu)\ge 0$, with equality if and only if
$\beta_3=0$ and $\beta_4=1$, i.e.\ $\nu=\tfrac12(\delta_{\bar\gamma+\sigma}+\delta_{\bar\gamma-\sigma})$.
\end{lemma}

\begin{proof}
By definition $\Delta_6(\nu):=-\tfrac{7}{64}\sigma^6-C_6$.
Using the explicit formula for $C_6$ from
Theorem~\ref{thm:poisson-kl-eighth},
\[
-\frac{7}{64}\sigma^6-C_6
=-\frac{7}{64}\sigma^6-\frac{\sigma^6+6\mu_3^2-8\sigma^2\mu_4}{64}
=\frac{-7\sigma^6-\sigma^6-6\mu_3^2+8\sigma^2\mu_4}{64}
=\frac{8\sigma^2\mu_4-8\sigma^6-6\mu_3^2}{64}.
\]
Writing $\mu_3=\sigma^3\beta_3$ and $\mu_4=\sigma^4\beta_4$ gives
\[
\Delta_6(\nu)
=\frac{8\sigma^6\beta_4-8\sigma^6-6\sigma^6\beta_3^2}{64}
=\frac{\sigma^6}{8}(\beta_4-1)-\frac{6\sigma^6\beta_3^2}{64}
=\frac{\sigma^6}{8}(\beta_4-1-\tfrac{3}{4}\beta_3^2).
\]
Since $\beta_4-1-\beta_3^2=\|Q_2\|^2\ge 0$ (the Pearson inequality)
and $\beta_4-1-\tfrac{3}{4}\beta_3^2\ge 0$ follows from the same bound,
one obtains the formula
\[
\Delta_6(\nu)
=\frac{\sigma^6}{32}\beta_3^2+\frac{\sigma^6}{8}(\beta_4-1-\beta_3^2),
\]
which is \eqref{eq:delta6-fourth-moment-lemma}.
(The split is verified by direct arithmetic:
$\sigma^6\beta_3^2/32+\sigma^6(\beta_4-1-\beta_3^2)/8
=\sigma^6\beta_3^2/32+\sigma^6\beta_4/8-\sigma^6/8-\sigma^6\beta_3^2/8
=\sigma^6(\beta_4-1)/8-\sigma^6\cdot 3\beta_3^2/32$,
matching the expression above.)
Non-negativity follows from $\beta_3^2\ge 0$ and $\beta_4-1-\beta_3^2\ge 0$.
Equality requires both $\beta_3=0$ and $\beta_4=1$, hence $\EE[X^2]=1$
and $\EE[X^4]=1$, which forces $X^2=1$ a.s., i.e.\ $\nu$ is the symmetric
two-point law.
\end{proof}

\begin{theorem}[Quantitative rigidity towards the symmetric two-point law]%
\label{thm:poisson-defect-stability}
Let $\nu$ be a probability measure on~$\RR$ with mean~$\bar\gamma$,
variance $\sigma^2>0$, and finite centred fourth moment.  Let
$\nu_c$ be the centred law of $\Delta=\gamma-\bar\gamma$,
$X:=\Delta/\sigma$, $\beta_3:=\EE[X^3]$, $\beta_4:=\EE[X^4]$,
and define
\begin{equation}\label{eq:delta6-fourth-moment}
\Delta_6(\nu)
:=\frac{\sigma^6}{32}\beta_3^2
+\frac{\sigma^6}{8}(\beta_4-1-\beta_3^2)
\end{equation}
as in Lemma~\ref{lem:delta6-fourth-moment}.
Set
\[
\rho_\sigma:=\frac12(\delta_{-\sigma}+\delta_\sigma).
\]
Then
\[
W_2(\nu_c,\rho_\sigma)
\le
2^{7/4}\sigma^{-1/2}\Delta_6(\nu)^{1/4}
+4\sqrt2\,\sigma^{-2}\Delta_6(\nu)^{1/2}.
\]
In particular, for fixed $\sigma$,
\[
W_2(\nu_c,\rho_\sigma)=O\!\bigl(\Delta_6(\nu)^{1/4}\bigr)
\qquad\text{as }\Delta_6(\nu)\downarrow 0.
\]
\end{theorem}

\begin{proof}
Keep the normalized variable $X=(\gamma-\bar\gamma)/\sigma$, set
\[
Q_2:=X^2-1-\beta_3X,
\]
and define
\[
Y:=\mathbf{1}_{\{X\ge 0\}}-\mathbf{1}_{\{X<0\}}\in\{-1,1\}.
\]
For every real $x$ one has
\[
|x-\operatorname{sgn}(x)|\le |x^2-1|,
\]
with the convention $\operatorname{sgn}(0)=1$.  Hence
\[
\EE[(X-Y)^2]\le \EE[(X^2-1)^2].
\]
Since $X^2-1=\beta_3X+Q_2$ and
\[
\EE[XQ_2]=\EE[X^3]-\EE[X]-\beta_3\EE[X^2]=0,
\]
one has the exact orthogonal decomposition
\[
\EE[(X^2-1)^2]
=\beta_3^2+\|Q_2\|_{L^2(\nu)}^2.
\]
The formula for $\Delta_6(\nu)$ from
Lemma~\ref{lem:delta6-fourth-moment} therefore yields
\[
\EE[(X-Y)^2]
\le \beta_3^2+\|Q_2\|_{L^2(\nu)}^2
\le 4\left(\frac{\beta_3^2}{4}+\|Q_2\|_{L^2(\nu)}^2\right)
=\frac{32}{\sigma^6}\Delta_6(\nu).
\]
Since $\EE[X]=0$,
\[
|\EE[Y]|=|\EE[Y-X]|
\le \EE[|Y-X|]
\le \EE[(Y-X)^2]^{1/2}
\le \frac{4\sqrt2}{\sigma^3}\Delta_6(\nu)^{1/2}.
\]
Let $Z$ be a symmetric Rademacher random variable.  Because both
$\mathcal L(Y)$ and $\mathcal L(Z)$ are supported on $\{-1,1\}$,
\[
W_2(\mathcal L(Y),\mathcal L(Z))^2=2|\EE[Y]|.
\]
Thus
\[
W_2(\mathcal L(Y),\mathcal L(Z))
\le \frac{2^{7/4}}{\sigma^{3/2}}\Delta_6(\nu)^{1/4}.
\]
Also,
\[
W_2(\mathcal L(X),\mathcal L(Y))
\le \EE[(X-Y)^2]^{1/2}
\le \frac{4\sqrt2}{\sigma^3}\Delta_6(\nu)^{1/2}.
\]
Since $\nu_c=\mathcal L(\sigma X)$ and $\rho_\sigma=\mathcal L(\sigma Z)$,
the scaling property of $W_2$ and the triangle inequality give
\begin{align*}
W_2(\nu_c,\rho_\sigma)
&=\sigma W_2(\mathcal L(X),\mathcal L(Z)) \\
&\le \sigma W_2(\mathcal L(X),\mathcal L(Y))
+\sigma W_2(\mathcal L(Y),\mathcal L(Z)) \\
&\le
4\sqrt2\,\sigma^{-2}\Delta_6(\nu)^{1/2}
+2^{7/4}\sigma^{-1/2}\Delta_6(\nu)^{1/4},
\end{align*}
which is the stated estimate.
\end{proof}

\begin{remark}\label{rem:w2-stability-constants}
The proof above is still coarse and we do not claim sharp constants in
Theorem~\ref{thm:poisson-defect-stability}.  The displayed estimate is
already stronger than the rough bound obtained from
$(a+b)^2\le 2a^2+2b^2$, because the identity $\EE[XQ_2]=0$ gives the
exact decomposition
$\EE[(X^2-1)^2]=\beta_3^2+\|Q_2\|_{L^2(\nu)}^2$.
\end{remark}

\subsection{A worked asymmetric three-point example}

\begin{example}[An asymmetric three-point law]
\label{ex:asymmetric-three-point}
Consider the centred probability measure
\[
\nu_{\mathrm{asym}}
:=\frac12\delta_{-1}+\frac14\delta_{0}+\frac14\delta_{2}.
\]
Then $\bar\gamma=0$,
\[
\sigma^2=\frac32,
\qquad
\mu_3=\frac32,
\qquad
\mu_4=\frac92,
\qquad
\mu_5=\frac{15}{2},
\qquad
\mu_6=\frac{33}{2}.
\]
Substituting these moments into
Theorems~\ref{thm:poisson-kl-eighth},
\ref{thm:poisson-defect-ladder}, and
\ref{thm:poisson-defect-amplification} gives
\[
D_{\mathrm{KL}}(P_t*\nu_{\mathrm{asym}}\|P_t)
=\frac{9}{32t^4}-\frac{297}{512t^6}+\frac{4617}{4096t^8}+o(t^{-8}),
\]
\[
\Delta_6(\nu_{\mathrm{asym}})=\frac{27}{128},
\qquad
\Delta_8(\nu_{\mathrm{asym}})=\frac{1377}{2048},
\]
and
\[
\frac{13\sigma^2}{8}\Delta_6(\nu_{\mathrm{asym}})
=\frac{1053}{2048},
\qquad
\Delta_8(\nu_{\mathrm{asym}})
-\frac{13\sigma^2}{8}\Delta_6(\nu_{\mathrm{asym}})
=\frac{81}{512}.
\]
Thus $\nu_{\mathrm{asym}}$ is genuinely asymmetric
($\mu_3,\mu_5\neq 0$), has a positive sixth-order Pearson gap, and
exhibits strict eighth-order defect amplification.

\begin{center}
\small
\begin{tabular}{lll}
\textbf{Quantity} & \textbf{Exact value} & \textbf{Role} \\
$C_6$ & $-\dfrac{297}{512}$ & sixth-order entropy coefficient \\
$\Delta_6(\nu_{\mathrm{asym}})$ & $\dfrac{27}{128}$ & Pearson-level defect \\
$C_8$ & $\dfrac{4617}{4096}$ & eighth-order entropy coefficient \\
$\Delta_8(\nu_{\mathrm{asym}})$ & $\dfrac{1377}{2048}$ & eighth-order defect \\
$\Delta_8-\dfrac{13\sigma^2}{8}\Delta_6$ & $\dfrac{81}{512}$ & strict amplification gap
\end{tabular}
\end{center}
\end{example}

\begin{lemma}[Kernel control of the fluctuation profile]\label{lem:rkhs-delta}
Assume $\nu$ has finite first absolute moment and let
\[
\widetilde\delta_t(y)
:=\frac{h_t(\bar\gamma+ty)}{g_t(\bar\gamma+ty)}-1
=\int_\RR \bigl(R_{\Delta/t}(y)-1\bigr)\,\dd\nu(\gamma),
\qquad \Delta:=\gamma-\bar\gamma.
\]
Then $\widetilde\delta_t\in L^2_0(\omega)$ and, under the unitary
$\mathcal U$ from Proposition~\ref{thm:mode-space-rkhs},
\[
(\mathcal U\widetilde\delta_t)(a)
=\frac{1}{t}\int_\RR \Delta\,K\!\left(a,\frac{\Delta}{t}\right)\dd\nu(\gamma),
\qquad a\in\RR.
\]
Moreover,
\[
\|\mathcal U\widetilde\delta_t\|_{\mathcal H_K}
=\|\widetilde\delta_t\|_{L^2(\omega)}
\le \frac{\EE[|\Delta|]}{\sqrt{2}\,t}.
\]
\end{lemma}

\begin{proof}
Since $R_{\varepsilon}(y)-1=\varepsilon\Psi_\varepsilon(y)$,
\[
\widetilde\delta_t(y)
=\frac{1}{t}\int_\RR \Delta\,\Psi_{\Delta/t}(y)\dd\nu(\gamma),
\]
which lies in $L^2_0(\omega)$ because each $\Psi_{\varepsilon}$ does.
Proposition~\ref{thm:mode-space-rkhs} gives
$\mathcal U(\Psi_\varepsilon)=K(\cdot,\varepsilon)$, hence the stated
Bochner integral formula for $(\mathcal U\widetilde\delta_t)(a)$.
The norm bound follows from
$\|\Psi_\varepsilon\|_{L^2(\omega)}^2=K(\varepsilon,\varepsilon)=1/2$
and the triangle inequality.
\end{proof}

\begin{corollary}[Strip control of the fluctuation profile]
\label{cor:rkhs-strip-control}
Under the assumptions of Lemma~\ref{lem:rkhs-delta}, the function
$\mathcal U\widetilde\delta_t$ extends holomorphically to the strip
\[
\mathbb S:=\{z\in\CC:|\Im z|<1\},
\]
and for every $a\in\RR$ and every $b\in(-1,1)$ one has
\[
\bigl|(\mathcal U\widetilde\delta_t)(a+\ii b)\bigr|
\le \frac{\EE[|\Delta|]}{2t\sqrt{1-b^2}}.
\]
\end{corollary}

\begin{proof}
By Lemma~\ref{lem:rkhs-delta},
$\mathcal U\widetilde\delta_t\in\mathcal H_K$ and
\[
\|\mathcal U\widetilde\delta_t\|_{\mathcal H_K}
\le \frac{\EE[|\Delta|]}{\sqrt2\,t}.
\]
Proposition~\ref{thm:gram-space-poisson-image} gives the holomorphic
extension to~$\mathbb S$ together with the pointwise bound
\[
\bigl|(\mathcal U\widetilde\delta_t)(a+\ii b)\bigr|
\le \frac{\|\mathcal U\widetilde\delta_t\|_{\mathcal H_K}}
{\sqrt{2(1-b^2)}},
\]
which yields the stated estimate.
\end{proof}

\subsection{Auxiliary kernel material}

The results above---the eighth-order expansion, the defect-ladder
decomposition, the amplification inequality, and the $W_2$-stability
theorem---form the main contribution of the paper.
Lemma~\ref{lem:rkhs-delta} and
Corollary~\ref{cor:rkhs-strip-control} show that the centred
fluctuation profile produced in Section~\ref{sec:precision-potential}
lives in the RKHS of Appendix~\ref{app:gram-space} and satisfies an
explicit strip bound there.  In that sense Appendix~\ref{app:gram-space}
feeds back into the same fluctuation profile that underlies the entropy
expansion.
Appendix~\ref{app:gram-space} computes the Gram kernel
$\langle\Psi_\varepsilon,\Psi_\eta\rangle_{L^2(\omega)}=2/(4+(\varepsilon-\eta)^2)$
and identifies the associated RKHS with the Poisson image space
$\sqrt\pi\,T_1(L^2(\RR))$.
Appendix~\ref{app:strip-lattice} records the explicit lattice symbol
and cardinal norm formula for the integer secant orbit; it is included
as an auxiliary closed-form specialization, and the structural results
are direct consequences of
\cite{deBoorDeVoreRon1994,RonShen1995,AldroubiGrochenig2001}.


